\section{Formulação de Elementos Finitos para Escoamentos em Meios Porosos}

Nesta seção, apresenta-se de modo sucinto a formulação clássica do problema de escoamento em meios porosos, bem como a formulação baseada na abordagem \textit{data-driven computing}. O conteúdo aqui presente, muito embora apresentado de forma enxuta, delimita os principais conceitos a serem empregados no desenvolvimento da nova técnica e elucidam a metodologia da problemática. 

\subsection{Formulação clássica do problema de escoamento em meios porosos}

\subsubsection{Formulação variacional contínua}

A lei constitutiva de Darcy, como fora referida outrora, advém do estudo realizado pelo físico francês Hanry Darcy, em 1856. Em seu trabalho, Hanry Darcy investigou os sistemas hidrológicos de suprimento de água na cidade de Dijon, na França. Para tanto, performou uma análise sobre o fluxo de água unidimensional em colunas de areia uniforme \cite{Darcy1856}. A partir do experimento, Darcy pôde concluir que o fluxo volumétrico, $Q$ (volume por unidade de tempo), era diretamente proporcional à área de seção transversal, $A$, à diferença $h_{t}-h_{b}$ e inversamente proporcional à altura da coluna de areia, $L$. A constante de proporcionalidade entre essas gradezas, $k$, foi denominada por Darcy de condutividade hidráulica.

Para o caso particular, a condutividade hidráulica é definida como $k = \frac{\rho g K}{\mu}$, onde $\mu$ denomina a viscosidade dinâmica do fluido, $ \rho $ a densidade do fluido, $ g $ a aceleração gravitacional e $ \mathbf{K} $ o tensor de permeabilidade absoluta\footnote{Se o meio é dito isotrópico, $ \mathbf{K} $ reduz-se a um escalar.}. Alguns autores, \textit{e.g.} \cite{Sousa2023}, denominam a razão $\frac{Q}{A}$ como sendo o fluxo de Darcy, $\mathbf{q}$; outros -- \textit{e.g.} \cite{Jambhekar2011} --, contudo, optam por denominar velocidade de Darcy, $\mathbf{u}$; neste último caso, no entanto, deve-se salientar que como apenas uma fração da área da seção transversal é considerada para o fluxo, em detrimento da presença dos grãos de areia, $\mathbf{u}$ corresponde a uma velocidade macroscópica média dos fluxos microscópicos que ocorrem no interior de REVs \cite{Sousa2023,Freeze1979}.     
Em coordenadas cartesianas e definindo $\mathbf{\hat{e}}_{z}$ vertical para cima, a lei de Darcy pode ser reescrita na forma:

\begin{equation}\label{eq5.1}
	\mathbf{u} = - \frac{\mathbf{K}}{\mu} \cdot (\nabla p - \rho \mathbf{g}),
\end{equation}

onde $\mathbf{K}$ é um tensor simétrico, uniformemente definido positivo definido com componentes em $L^{\infty}(\Omega)$  e $\mathbf{g} = -g\mathbf{\hat{e}}_{z}$. 

Para estabelecer as equações diferenciais que modelam o fluxo em meios permeáveis, combinadas com a lei de Darcy e as condições de contorno, são empregados os princípios de conservação de massa e de energia. No entanto, aqui nos restringimos a fenômenos isotérmicos, portanto, utilizaremos apenas o princípio de conservação de massa, cuja equação é dada por:

\begin{equation}\label{eq5.2}
	\dfrac{\partial }{\partial t}(\phi \rho) + \nabla \cdot (\rho \mathbf{u}) = f,
\end{equation}

onde $\phi$ é a porosidade  da rocha e $f \in L^{2}(\Omega)$ delimita um termo fonte (ou sorvedouro). As equações \ref{eq5.1} e \ref{eq5.2}, conjuntamente com as condições de contorno definem o problema denominado de Darcy para escoamento monofásico. Assim sendo, seja $\Omega \subset \mathbb{R}^{d}$ o domínio sobre o qual o problema é definido. Admitindo-se escoamento incompressível e em regime permanente, e desprezando-se as variações gravitacionais, as equações \ref{eq5.1} e \ref{eq5.2} se reduzem a:

\begin{equation}
	\begin{cases}
		\mathbf{u} = -\boldsymbol{\kappa} \cdot \nabla p, \quad \boldsymbol{\kappa} = \frac{\mathbf{K}}{\mu} \ \mid \ \boldsymbol{\kappa} \in [L^{\infty}(\Omega)]^{d \times d} \\
		\nabla \cdot \mathbf{u} = f
	\end{cases}
\end{equation}

Nesse sentido, introduzindo-se as condições de contorno ao problema, tanto para o campo escalar de pressões, $p$, quanto para o fluxo, $\mathbf{u}$, resulta que:

\begin{equation}\label{eq5.4}
	\begin{cases}
		\mathbf{u}= -\boldsymbol{\kappa} \cdot \nabla{p} & \text{em}~~\Omega \\
		\nabla \cdot \mathbf{u} = f & \text{em}~~\Omega \\
		p = p_D & \text{em}~~\Gamma_D \\
		\mathbf{u} \cdot \hat{\mathbf{n}} = u_N & \text{em}~~\Gamma_N \\
	\end{cases}
\end{equation}

onde $\Gamma_{D}$ refere-se à fronteira do domínio sobre a qual foi prescrita uma condição de contorno de Dirichlet e, por outro lado, $\Gamma_{N}$ refere-se à fronteira do domínio sobre a qual foi imposta uma condição de contorno de Neumann. Notoriamente, $\partial \Omega = \bigcup_{i} \Gamma_{i}$ e $\bigcap_{i} \Gamma_{i} = \emptyset$. Multiplicando por $\boldsymbol{\kappa}^{-1}$ ambos os lados da primeira equação em \ref{eq5.4} pode-se definir o problema variacional ou a formulação fraca equivalente, baseada na formulação de Garlerkin, como sendo: Encontrar $(\mathbf{u}, p) \in V_{u} \times Q$, com $V_{u} = \{ \mathbf{v} \in \text{Hdiv}(\Omega) \mid \mathbf{v} \cdot \mathbf{\hat{n}} = u_{N}, \ \forall \mathbf{x} \in \Gamma_{N}\}$, tais que

\begin{equation}
	\begin{aligned}
		\int_{\Omega}\boldsymbol{\kappa}^{-1} \cdot \mathbf{u} \cdot \mathbf{v} \, d\mathbf{x}  - \int_{\Omega}p\nabla \cdot \mathbf{v} \, d\mathbf{x} &= -\int_{\Gamma_{D}}p_D \left (\mathbf{v} \cdot \hat{\mathbf{n}}\right ) \, d\mathbf{s} \\
		\int_{\Omega} q \nabla \cdot \mathbf{u} d\mathbf{x} &= \int_{\Omega} f \, q \,d\mathbf{x}
	\end{aligned}
\end{equation}

$\forall (\mathbf{v},q) \in V_0 \times Q$, onde $V_{0} = \{ \mathbf{v} \in \text{Hdiv}(\Omega)  \mid \mathbf{v} \cdot \mathbf{\hat{n}} = 0, \ \forall \mathbf{x} \in \partial \Omega\}$ e $Q = L^{2}(\Omega) = \{q:\Omega \rightarrow\mathbb{R} \mid \int_{\Omega} |q(x)|^{2} d\Omega < \infty\}$.

Note que a formulação do problema de Darcy é equivalente à formulação mista do problema de Poisson. Outrossim, a condição de Neumann sobre $\mathbf{u}$ em $\Gamma_N$ surge de forma forte como uma restrição no conjunto de testes $V_{u}$ para $\mathbf{u}$, enquanto a condição de Dirichlet para $p_D$ em $\Gamma_D$ surge de forma natural no problema variacional \cite{Ausas2022}.

Um caso particular do problema definido pela equação \ref{eq5.4} consiste naquele em que apenas condições de contorno do tipo Dirichlet são impostas na fronteira do domínio. Nesse cenário, o problema \ref{eq5.4} se reduz a:

\begin{equation}\label{eq5.6}
	\begin{cases}
		\mathbf{u}= -\boldsymbol{\kappa} \cdot \nabla{p} & \text{em}~~\Omega \\
		\nabla \cdot \mathbf{u} = f & \text{em}~~\Omega \\
		p = p_D & \text{em}~~\partial \Omega
	\end{cases}
\end{equation}

A formulação variacional, com isso, torna-se: Encontrar $(\mathbf{u},p) \in V \times Q$, tais que

\begin{equation}
	\begin{aligned}
		\int_{\Omega}\boldsymbol{\kappa}^{-1} \cdot \mathbf{u} \cdot \mathbf{v} \, d\mathbf{x}  - \int_{\Omega}p\nabla \cdot \mathbf{v} \, d\mathbf{x} &= -\int_{\partial \Omega}p_D \left (\mathbf{v} \cdot \hat{\mathbf{n}}\right ) \, d\mathbf{s} \\
		\int_{\Omega} q \nabla \cdot \mathbf{u} d\mathbf{x} &= \int_{\Omega} f \, q \,d\mathbf{x}
	\end{aligned}
\end{equation}

$\forall (\mathbf{v},q) \in V_0 \times Q$.

Por fim, cabe a ressalva de que na modelagem computacional do fluxo em meios porosos, a seleção adequada de subespaços discretos dentro do espaço de Hilbert $\text{Hdiv}(\Omega)$ -- i.e. $\text{Hdiv}(\Omega) = \{\mathbf{v} \in [L^{2}(\Omega)]^{d} \mid \nabla \cdot \mathbf{v} \in L^{2}(\Omega)\}$ -- é fundamental para assegurar a conservação local, de forma análoga aos métodos de volumes finitos e os métodos de Galerkin descontínuo.

\subsubsection{Formulação variacional discreta}

A formulação variacional discreta do problema é obtida considerando subespaços de dimensão finita dos espaços em que foi formulado o problema contínuo, ou seja:

\begin{equation}
	V_{h,u} \subset V_u \subset \text{Hdiv}(\Omega), \quad Q_h \subset L^2(\Omega)
\end{equation}

O problema variacional discreto é dado por: Encontrar \((\mathbf{u}_h , p_h) \in V_{h,u} \times Q_h\), tais que

\begin{equation}
	\begin{aligned}
		a(\mathbf{u}_h , \mathbf{v}_h) + b(\mathbf{v}_h , p_h) = c(\mathbf{v}_h)\\
		b(\mathbf{u}_h , q_h) = \ell(q_h)
	\end{aligned}
\end{equation}
para todo \((\mathbf{u}_h , q_h) \in V_{h,0} \times Q_h\),  com as formas bilineares e lineares dadas por:
\begin{equation}
	\begin{aligned}
		&a(\mathbf{u}_h, \mathbf{v}_{h}) = \int_{\Omega} \boldsymbol{\kappa}^{-1} \cdot \mathbf{u}_{h} \cdot \mathbf{v}_{h} d\mathbf{x}, \quad b(\mathbf{v}_{h},p_{h}) = -\int_{\Omega} p_{h} \nabla \cdot \mathbf{v}_{h} \, d\mathbf{x}, \\
		&c(\mathbf{v}_{h}) = -\int_{\partial \Omega} p_{D} \mathbf{v}_{h} \cdot \mathbf{\hat{n}} \, d\mathbf{s}, \quad b(\mathbf{u}_{h},q_{h}) = -\int_{\Omega} q_{h} \nabla \cdot \mathbf{u}_{h} \, d\mathbf{x}\\
		&\ell(q) = -\int_{\Omega} f q \, d\mathbf{x}.
	\end{aligned}
\end{equation}


\subsection{Formulação \textit{data-driven computing} do problema de escoamento em meios porosos}

Uma vez definida a formulação clássica do problema de Darcy, apresenta-se aqui a formulação baseada na abordagem \textit{data-driven computing} (c.f. seção \ref{sec:introducao}). Assim, denotemos de espaço de fase, $Z$, o espaço que contém os pares $(\mathbf{u}_{d},\mathbf{g}_{d})$ -- i.e. fluxo e gradiente de pressão, respectivamente -- que definem o estado local do escoamento no meio poroso. Cada ponto no conjunto de dados corresponde, por exemplo, a uma medida experimental. Ademais, seja $\Omega \subset \mathbb{R}^{d}$ um domínio limitado e simplesmente conexo sobre o qual define-se um estado por meio de um campo escalar de pressões $p:\Omega \rightarrow \mathbb{R}$. Com isso, pode-se definir o problema equivalente à equação \ref{eq5.6}, sob a ótica da abordagem \textit{data-driven computing}, como sendo: Encontrar $\mathbf{u},p,\mathbf{g}, \lambda, \boldsymbol{\mu}$ que minimizam o funcional $\boldsymbol{\mathcal{J}}$, definido por:

\begin{equation}
	\boldsymbol{\mathcal{J}} = \dfrac{1}{2} \int_{\Omega} \ell \big( (\mathbf{u} - \mathbf{u}_{\text{d}}). (\mathbf{g} - \mathbf{g}_{\text{d}})\big)d\mathbf{x},
\end{equation}

com $ (\mathbf{u}_{\text{d}},\textbf{g}_{\text{d}}) $ dados e sujeitos as seguintes restrições:

\begin{equation}
	\begin{aligned}
		\textbf{g} -  \nabla p &= 0, \quad &\text{em} \ \Omega \\
		\nabla \cdot \textbf{u} - f &= 0, \quad &\text{em} \ \Omega \\
		p - p_{D} &= 0, \quad  &\text{em} \ \partial \Omega \\
	\end{aligned}
\end{equation}

Define-se, então, o Lagrangiano como sendo:

\begin{equation}\label{eq5.10}
	\boldsymbol{\mathcal{L}} = \boldsymbol{\mathcal{J}} + \int_{\Omega} \lambda \big(\nabla \cdot \textbf{u} - f\big)d\mathbf{x} + \int_{\Omega} \boldsymbol{\mu} \cdot (\textbf{g} -  \nabla p)d\mathbf{x},
\end{equation}

onde $\lambda$ e $\boldsymbol{\mu}$ denominam os multiplicadores de Lagrange para as equações de balanço e de compatibilidade, respectivamente. Deve-se observar que, a depender da escolha do espaço de elementos finitos que se deseja empregar, diferentes conjuntos de condições de optimalidade tornam-se possíveis. Na primeira abordagem, a qual denotaremos por 5A, integra-se por partes a restrição de \textbf{compatibilidade}, ou seja:

\begin{equation}
	\begin{aligned}
		\int_{\Omega} \boldsymbol{\mu} \cdot (\mathbf{g} -  \nabla p)d\mathbf{x} &= \int_{\Omega}(\boldsymbol{\mu} \cdot \mathbf{g})d\mathbf{x} - \int_{\Omega}(\boldsymbol{\mu} \cdot  \nabla p)d\mathbf{x} \\
		& = \int_{\Omega}(\boldsymbol{\mu} \cdot \mathbf{g})d\mathbf{x} - \int_{\Omega}\nabla \cdot (p\boldsymbol{\mu}) + \int_{\Omega}p\nabla \cdot \boldsymbol{\mu}d\mathbf{x} \\
		& = \int_{\Omega}(\boldsymbol{\mu} \cdot \textbf{g} )d\mathbf{x} + \int_{\Omega}p\nabla \cdot \boldsymbol{\mu}d\mathbf{x} - \int_{\partial \Omega} p_{D} \boldsymbol{\mu} \cdot  \mathbf{\hat{n}}d\mathbf{s} 
	\end{aligned}
\end{equation}

Com isso, o Lagrangiano assume a seguinte forma:

\begin{equation}
	\boldsymbol{\mathcal{L}} = \boldsymbol{\mathcal{J}} + \int_{\Omega} \lambda \big(\nabla \cdot \mathbf{u} - f\big)dv +\int_{\Omega}(\boldsymbol{\mu} \cdot \textbf{g} + p\nabla \cdot \boldsymbol{\mu})d\mathbf{x} - \int_{\partial \Omega} p_{D} \boldsymbol{\mu} \cdot  \mathbf{\hat{n}}d\mathbf{s}.
\end{equation}

Tomando-se, então, todas as possíveis variações, obtém-se que:

\begin{equation}
	\begin{aligned}
		&d_{\mathbf{u}}\boldsymbol{\mathcal{L}} = d_{\mathbf{u}}\boldsymbol{\mathcal{J}} + \int_{\Omega}(\lambda\nabla \cdot \delta \mathbf{u})d\mathbf{x}, &\quad \forall \delta \mathbf{u} \\
		&d_{p}\boldsymbol{\mathcal{L}} = \int_{\Omega}(\delta p \nabla \cdot \boldsymbol{\mu}) d\mathbf{x}, &\quad \forall \delta p \\
		&d_{\mathbf{g}}\boldsymbol{\mathcal{L}} = d_{\mathbf{g}}\boldsymbol{\mathcal{J}}  + \int_{\Omega} (\delta\mathbf{g} \cdot \boldsymbol{\mu})d\mathbf{x}, &\quad \forall \delta \mathbf{g} \\
		&d_{\lambda}\boldsymbol{\mathcal{L}} = \int_{\Omega} \delta \lambda (\nabla \cdot \mathbf{u} - f)d\mathbf{x}, &\quad \forall \delta \lambda \\
		&d_{\boldsymbol{\mu}}\boldsymbol{\mathcal{L}} = \int_{\Omega}(\mathbf{g} \cdot \delta \boldsymbol{\mu} + p \nabla \cdot \delta \boldsymbol{\mu})d\mathbf{x} - \int_{\partial \Omega} (p_{D} \delta \boldsymbol{\mu} \cdot  \mathbf{\hat{n}}) d\mathbf{s}, &\quad \forall \delta \boldsymbol{\mu}
	\end{aligned}
\end{equation}

Noutra vertente, na segunda formulação, a qual atribuiremos a nomenclatura 5B, integra-se por partes a equação de \textbf{balanço}, ou seja:

\begin{equation}
	\begin{aligned}
		\int_{\Omega}\lambda (\nabla \cdot \mathbf{u} - f)d\mathbf{x} & = \int_{\Omega}\lambda \nabla \cdot \mathbf{u} d\mathbf{x} - \int_{\Omega} \lambda f d\mathbf{x} \\		
		& = \int_{\Omega} \nabla \cdot (\lambda \mathbf{u}) d\mathbf{x} - \int_{\Omega} \mathbf{u} \cdot \nabla \lambda d\mathbf{x} - \int_{\Omega} \lambda f d\mathbf{x} \\
		& = \underbrace{\int_{\partial \Omega} \lambda \mathbf{u} \cdot \mathbf{\hat{n}} d\mathbf{s}}_{\lambda|_{\partial \Omega} = 0 } - \int_{\Omega} \mathbf{u} \cdot \nabla \lambda d\mathbf{x} - \int_{\Omega} \lambda f d\mathbf{x} \\
		& = - \int_{\Omega} \mathbf{u} \cdot \nabla \lambda d\mathbf{x} - \int_{\Omega} \lambda f d\mathbf{x}.
	\end{aligned}
\end{equation}

Desta forma, o Lagrangiano toma a seguinte forma: 

\begin{equation}
	\boldsymbol{\mathcal{L}} = \boldsymbol{\mathcal{J}}  - \int_{\Omega} \mathbf{u} \cdot \nabla \lambda d\mathbf{x} - \int_{\Omega} \lambda f d\mathbf{x} + \int_{\Omega} \boldsymbol{\mu} \cdot (\textbf{g} -  \nabla p)d\mathbf{x}
\end{equation}

Tomando-se, então, todas as possíveis variações, resulta que:

\begin{equation}
	\begin{aligned}
		&d_{\mathbf{u}}\boldsymbol{\mathcal{L}} = d_{\mathbf{u}}\boldsymbol{\mathcal{J}} - \int_{\Omega}(\delta \mathbf{u} \cdot \nabla \lambda )d\mathbf{x}, &\quad \forall \delta \mathbf{u} \\
		&d_{p}\boldsymbol{\mathcal{L}} =  - \int_{\Omega}( \boldsymbol{\mu} \cdot  \nabla \delta p )d\mathbf{x}, &\quad \forall \delta p \\
		&d_{\mathbf{g}}\boldsymbol{\mathcal{L}} = d_{\mathbf{g}}\boldsymbol{\mathcal{J}}  + \int_{\Omega} (\delta \mathbf{g} \cdot \boldsymbol{\mu})d\mathbf{x}, &\quad \forall \delta \mathbf{g} \\
		&d_{\lambda}\boldsymbol{\mathcal{L}} = -\int_{\Omega} \big(\mathbf{u} \cdot \nabla \delta \lambda + f \delta \lambda \big)d\mathbf{x}, &\quad \forall \delta \lambda \\
		&d_{\boldsymbol{\mu}}\boldsymbol{\mathcal{L}} = \int_{\Omega} \delta \boldsymbol{\mu} \cdot (\textbf{g} -  \nabla p)d\mathbf{x}, &\quad \forall \delta \boldsymbol{\mu}
	\end{aligned}
\end{equation}


Nesta altura, deve-se especificar a expressão para o funcional $\boldsymbol{\mathcal{J}}$, i.e. a função objetivo que almeja-se minimizar. Tomando-se, então, $\boldsymbol{\mathcal{J}} $ como sendo:

\begin{equation}
	\boldsymbol{\mathcal{J}} = \frac{1}{2} \int_{\Omega} \left[ (\mathbf{q} - \mathbf{q}_{\text{d}}) \cdot (\mathbf{q} - \mathbf{q}_{\text{d}}) + (\mathbf{g} - \mathbf{g}_{\text{d}}) \cdot (\mathbf{g} - \mathbf{g}_{\text{d}}) \right] d\mathbf{x},
\end{equation}

temos que a derivada direcional em relação a $\mathbf{q}$ é dada por:
\begin{equation}
	d_{\mathbf{q}}\boldsymbol{\mathcal{J}} = \lim_{\epsilon \to 0} \frac{\boldsymbol{\mathcal{J}}(\mathbf{q} + \epsilon \mathbf{\delta q}, \mathbf{g}) - \boldsymbol{\mathcal{J}}(\mathbf{q}, \mathbf{g})}{\epsilon}.
\end{equation}

Substituindo a expressão de $\boldsymbol{\mathcal{J}}$, obtemos que:
{\small
	\begin{equation}
		d_{\mathbf{q}}\boldsymbol{\mathcal{J}} = \lim_{\epsilon \to 0} \frac{1}{2\epsilon} \left[ \int_{\Omega} \left( (\mathbf{q} + \epsilon \mathbf{\delta q} - \mathbf{q}_{\text{d}}) \cdot (\mathbf{q} + \epsilon \mathbf{\delta q} - \mathbf{q}_{\text{d}}) - (\mathbf{q} - \mathbf{q}_{\text{d}}) \cdot (\mathbf{q} - \mathbf{q}_{\text{d}}) \right) d\mathbf{x} \right]
	\end{equation}
}

Expandindo o quadrado no integrando e cancelando termos semelhantes, obtemos:
\begin{equation}
	d_{\mathbf{q}}\boldsymbol{\mathcal{J}} = \lim_{\epsilon \to 0} \frac{1}{2\epsilon} \left[ \int_{\Omega} \left( 2\epsilon \mathbf{\delta q} \cdot (\mathbf{q} - \mathbf{q}_{\text{d}}) + \epsilon^2 \mathbf{\delta q} \cdot \mathbf{\delta q} \right) d\mathbf{x} \right]
\end{equation}
\begin{equation}
	d_{\mathbf{q}}\boldsymbol{\mathcal{J}} = \int_{\Omega} \mathbf{\delta q} \cdot (\mathbf{q} - \mathbf{q}_{\text{d}}) d\mathbf{x}
\end{equation}

Onde usamos o fato de que  $\lim_{\epsilon \to 0} \frac{\epsilon^2}{2\epsilon} = 0$. Para a derivada direcional em relação a $\mathbf{g}$, o procedimento é análogo, e obtemos:
\begin{equation}
	d_{\mathbf{g}}\boldsymbol{\mathcal{J}} = \int_{\Omega} \mathbf{\delta g} \cdot (\mathbf{g} - \mathbf{g}_{\text{d}}) d\mathbf{x}.
\end{equation}

Por fim, deve-se especificar a escolha dos espaços de elementos finitos para as variáveis do problema. Optamos, aqui, pelos requisitos mínimos de regularidade para os espaços de funções. Assim, considere uma partição $\mathcal{T}_h = \{K_i\}_{i=1}$ de $\Omega$, tal que $\overline{\Omega} \subseteq \bigcup_{i=1}^{N} \overline{K}_i$\footnote{O fecho (ou fecho topológico) do domínio \(\Omega\) está contido na união dos fechos dos elementos \(K_i\). Ou seja, a fronteira de \(\Omega\) está completamente contida na união das fronteiras dos elementos \(K_i\).}, onde $K_i \cap K_j = \emptyset$ se $i \neq j$. $\mathcal{T}_h$ é chamada de uma triangulação de $\Omega$. Dada essa triangulação, definiremos 

\begin{equation}
	\text{CG}_k(\mathcal{T}_h) = \{ v \in H^1(\Omega),~ v|_K \in P_k(K)~\forall K \in \mathcal{T}_h\},    
\end{equation}

como o espaço das funções globalmente continuas de grau $k$ com respeito a triangulação  $\mathcal{T}_{h}$ e $P_{k}(K)$ corresponde ao espaço das funções polinomiais de grau $k$ em cada elemento $K$ da triangulação $\mathcal{T}_{h}$. Adicionalmente, definiremos

\begin{equation}
	\text{DG}_k(\mathcal{T}_h) = \{  v \in L^2(\Omega),~ v|_K \in P_k(K)~\forall K \in \mathcal{T}_h \}
\end{equation}

como sendo o espaço das funções Galerkin descontinuas de grau $k$ com respeito a triangulação $\mathcal{T}_{h}$. Por fim, definimos:

\begin{equation}
	\text{Hdiv}_1 = \{ \mathbf{v} \in [L^2(\Omega)]^d \mid \nabla \cdot \mathbf{v} \in L^2(\Omega), \mathbf{v}|_K \in [P_k(K)]^d, \forall K \in \mathcal{T}_h \}.
\end{equation}

Isso posto, a tabela \ref{tab:espacos_elefim} resume a escolha dos espaços de elementos finitos para cada uma das variáveis do problema, onde $\text{CG}_k^{p}(\mathcal{T}_h)$ e $\text{CG}_k^{\lambda}(\mathcal{T}_h)$ são definidos, respectivamente, por:

\begin{equation}
	\text{CG}_k^{p}(\mathcal{T}_h) = \{ v \in H^1(\Omega),~ v|_K \in P_k(K)~\forall K \in \mathcal{T}_h,~v = p_D~\mbox{em}~\partial{\Omega} \}
\end{equation}

e 

\begin{equation}
	\text{CG}_k^{\lambda}(\mathcal{T}_h) = \{ v \in H^1(\Omega),~ v|_K \in P_k(K)~\forall K \in \mathcal{T}_h,~v = 0~\mbox{em}~\partial{\Omega}\}.  
\end{equation} 

\begin{table}[h!]
	\centering
	\caption{Escolha dos espaços de elementos finitos para as formulações 5A e 5B.}
	\label{tab:espacos_elefim}
	\begin{tabular}{ccc}
		\toprule
		\textbf{Variável} & \textbf{5A} & \textbf{5B} \\
		\midrule
		$\mathbf{u}$          & $\text{Hdiv}_{k+1}$ & $[\text{DG}_{k}]^{d}$ \\
		$\boldsymbol{\mu}$  & $\text{Hdiv}_{k+1}$ & $[\text{DG}_{k}]^{d}$ \\
		$p$    & $\text{DG}_{k}$      & $\text{CG}_{k+1}^{p}$     \\
		$\lambda$    & $\text{DG}_{k}$      & $\text{CG}_{k+1}^{\lambda}$     \\
		$\mathbf{g}$    & $[\text{DG}_{k}]^{d}$ & $[\text{DG}_{k}]^{d}$ \\
		\bottomrule
	\end{tabular}
\end{table}

Além disso, é relevante mencionar uma observação adicional. Pode-se obter outras formulações de elementos finitos para o problema sob a ótica da abordagem \textit{data-driven computing}. Ao integrarmos por partes ambas as equações - a de compatibilidade e a de balanço - obtemos uma nova formulação fraca, a qual podemos denominar de 5C. Uma quarta possibilidade, por sua vez, seria não integrar por partes nenhuma das equações, resultando em uma nova formulação, como por exemplo, a 5D.



